% Provisional abstract. The paragraph marked below is a placeholder until
% the CRESCO8 measurements exist; everything else reflects the background
% and methodology as they currently stand.
\begin{abstract}
    Running large language models on CPUs is no longer unusual, and the
    work that made it ordinary has been overwhelmingly arithmetic: fewer
    bits per weight, cheaper kernels, better use of the matrix units.
    That line of attack is visibly saturating. A recent ternary-kernel
    design reports a twenty-nine-fold improvement in isolation and a
    1.24-fold improvement end to end, and the distance between those two
    numbers is not an engineering failure but a statement about where the
    time now goes. Once a generation step spends most of its duration
    waiting for weights to arrive, removing instructions from it buys
    very little, and the interesting question becomes where those weights
    are kept.

    This thesis takes that question to hardware that can answer it
    directly. The Intel Xeon CPU Max~9480 exposes 64\,GB of on-package
    HBM2e alongside conventional DDR5 as separately addressable memory
    tiers on a single socket, so choosing between them is a page-placement
    decision rather than a data transfer, and it requires no change to the
    inference runtime. Using the HBM-equipped partition of the CRESCO8
    cluster at ENEA Portici, the study compares verified local-HBM against
    local-DDR5 placement for quantized inference under
    \texttt{llama.cpp}, holding the build, model, quantization format,
    thread count and prompt fixed and reporting prompt processing and
    sustained token generation as separate measurements rather than as a
    single throughput figure.

    Two predictions are registered in advance. Token generation should
    gain more from the faster tier than prompt processing, because its
    low-batch matrix--vector work sits far to the left of the Roofline
    ridge point where the two tiers' ceilings are furthest apart. The gain
    should nonetheless fall short of the ratio between the tiers' STREAM
    bandwidths, because sustaining HBM bandwidth requires roughly four
    times the memory-level parallelism that DDR5 does, and a dependent
    decode loop may not supply it. Both predictions can fail informatively:
    a small effect accompanied by low achieved bandwidth would locate the
    limit in concurrency rather than in the memory technology.

    The background developed here situates the problem between two
    literatures that have largely not met. CPU inference has concentrated
    on arithmetic; memory tiering has strong measured results in
    high-performance computing but has reached transformer inference
    mainly through simulation, through accelerator-attached hierarchies
    reached over an interconnect, or through designs that rebuilt the
    runtime around a different fast tier. A Xeon Max socket removes all
    three qualifications at once, which is what makes the measurement
    worth taking.

    % TODO: results and conclusions paragraph, after the CRESCO8
    % measurement campaign. State the measured phase-resolved ratios,
    % where the capacity boundary fell, and which hypothesis survived.
\end{abstract}
